\bigskip\Needspace{10\baselineskip}
\pdfbookmark[1]{Pallas Olbersiana. 1802 November}{pallas180211}
\begin{center}
PALLAS OLBERSIANA.\par
\medskip\hrule\medskip
Monatliche Correspondenz. Band VI. S.\ 499—505. 1802 November.\par
\medskip\hrule
\end{center}\medskip
—\quad —\quad —

Dr.\ \textsc{Herschel} war so freundlich, uns einen Sonderdruck seiner Abhandlung aus den \textit{Philosophical Transactions}\HistoricalNote{22} zu schicken. Ihren wesentlichen Inhalt haben wir unseren Lesern bereits im Juliheft auf S.\ 90 in einem Auszug aus einem seiner Briefe mitgeteilt. Nachdem wir die Abhandlung noch einmal aufmerksam gelesen haben, sehen wir keinen Anlass, unsere bereits geäußerte Meinung darüber zu ändern. Aus unserem Briefwechsel geht auch hervor, dass andere Astronomen, darunter \textsc{La Lande}, \textsc{Oriani}, \textsc{Piazzi} und \textsc{Gauß}, seiner neu vorgeschlagenen Klassifikation nicht zustimmen. Besondere Beachtung verdienen die Gründe und eine sehr treffende Bemerkung von Dr.\ \textsc{Gauß} gegen das angebliche sogenannte harmonische Gesetz, mit dem die Abstände der Planeten übereinstimmen sollen. Dieser scharfsinnige Mathematiker äußert sich dazu in einem Schreiben vom 16. Oktober auf bemerkenswerte Weise.

„Dr.\ \textsc{Herschel} will die neuen Planeten, wie mir Prof.\ \textsc{Huth} sagte, der ihn in England besucht hat, noch immer nicht dulden, obwohl meines Wissens noch kein einziger Astronom seinen Vorschlag gebilligt hat. Im Grunde hängt es, wie Sie so richtig bemerkt haben, nur von der Übereinkunft ab, ob wir \textit{Ceres} und \textit{Pallas} als Planeten ansehen wollen. Es geht dabei gar nicht darum, ob sie Planeten \textit{sind}, sondern darum, ob es \textit{angemessen} und \textit{passend} ist, diese Himmelskörper Planeten zu \textit{nennen}. Sie stimmen in einigen Punkten mit den \textit{bisher bekannten} Planeten überein und weichen in anderen von ihnen ab. Dass die letzteren Punkte ganz unwesentlich sind, haben Sie im Juliheft dieses Jahres zur Genüge gezeigt. (Bei \textit{Ceres} lässt sich die Bahnneigung, wie Sie im Dezemberheft des vergangenen Jahres bemerkt haben, nicht einmal zu diesen Punkten zählen.) Dass den Astronomen eine kreisähnliche Bahn und die davon abhängige \textit{dauernde Gegenwart} stets als das Wesentliche galten, scheint daraus hervorzugehen, dass
\SourceLeaf{236}%
die Astronomen sie sofort ohne Weiteres als Planeten anerkannten, sobald sie sich von dieser Beschaffenheit der Bahn überzeugt hielten. Ich glaube sogar, dass wir auch dann nicht unbedingt von der Bezeichnung als Planeten abgehen müssten, wenn die Zukunft die Vermutung unseres vortrefflichen \textsc{Olbers}, \textit{Ceres} und \textit{Pallas} seien nur Bruchstücke eines zerstörten Planeten, durch die Entdeckung weiterer Bruchstücke bestätigen sollte. Bei der Prüfung ihres Anspruchs auf die Bezeichnung als Planet scheint es mir nämlich stärker darauf anzukommen, \textit{ob} diese Himmelskörper ihrer wesentlichen Eigenschaft nach Planeten sind, als darauf, \textit{wie} sie es \textit{geworden} sind.

„Es ist sonderbar, dass man das von Prof.\ \textsc{Titius} angegebene sogenannte Gesetz als Argument gegen die beiden Planeten verwenden wollte. Anders als bei allen Wahrheiten, die den Namen Gesetz verdienen, trifft dieses Verhältnis bei den übrigen Planeten nur ganz zufällig zu. Beim Merkur trifft es überhaupt nicht zu, was man noch nicht einmal bemerkt zu haben scheint\hyperlink{native-236-note}{\textsuperscript{*)}}. Mir scheint sehr einleuchtend, dass die Folge
\[
4,\ 4+3,\ 4+6,\ 4+12,\ 4+24,\ 4+48,\ 4+96,\ 4+192
\]
mit der die Abstände übereinstimmen sollten, nicht einmal eine kontinuierliche Folge ist. Das Glied vor $4+3$ kann ja nicht $4$ sein, also $4+0$. Es muss $4+1\frac12$ sein. Demnach müssten zwischen $4$ und $4+3$ noch unendlich viele Glieder liegen. Oder, wie \textsc{Wurm} es ausdrückt (M.\ C., Juniheft 1801, S.\ 594), für Merkur oder für $n=1$ ergibt $4+2^{n-2}\,3$ nicht $4$, sondern $5\frac12$.

„Es ist gar nichts daran auszusetzen, wenn man solche ungefähren Übereinstimmungen in der Natur sucht. Die größten Männer haben sich solchem \textit{lusus ingenii} hingegeben. So viel sich \textsc{Kepler} auch auf seine regelmäßigen Körper zugutehielt, die er mit den Planetenabständen in Übereinstimmung gebracht hatte (er wollte, wie er sagte, die Ehre dieses Fundes nicht für das Kurfürstentum Sachsen hergeben), so hätte er doch, wenn Uranus zu seiner Zeit entdeckt worden wäre, diesem gewiss nicht deshalb abgesprochen, ein Planet zu sein, weil er nicht zu seinen Vorstellungen passte. Vielmehr hätte er diese ohne jeden Zweifel sofort aufgegeben. Ebenso hatten Sie es gemacht. Auch Sie haben über jenes Gesetz von \textsc{Titius} nachgedacht. Doch Sie haben Ihre Überlegungen bloße Träume genannt und sie sogleich aufgegeben, als sie anfingen, den Tatsachen zu widersprechen.“
\par\nobreak\smallskip
\noindent\rule{.24\linewidth}{.3pt}\par\nobreak\smallskip
{\footnotesize\indent\hypertarget{native-236-note}{}*)\enspace\csname HistoricalNoteText23\endcsname\par}
\begin{flushright}[\textsc{von Zach.}]\end{flushright}

